% Assessor pack, rendered as nine cards. Generated by scripts/gen-pack-cards.py.
% Each PNG is 2x so it stays crisp in print. Two per row at 0.48\textwidth.
\subsection*{A.x The assessor pack}

The run sheet the moderator worked from. It is reproduced as the nine cards it was used as,
rather than as continuous prose, because that is the form it took in the sessions: one card per
thing the moderator had to do or decide. Card 5 carries the success criteria and was deliberately
withheld from participants.

\begin{figure}[H]
\centering
\safeincludegraphics[width=0.48\textwidth]{pack-assessor-01.png}\hfill
\safeincludegraphics[width=0.48\textwidth]{pack-assessor-02.png}
\[0.7em]
\safeincludegraphics[width=0.48\textwidth]{pack-assessor-03.png}\hfill
\safeincludegraphics[width=0.48\textwidth]{pack-assessor-04.png}
\[0.7em]
\safeincludegraphics[width=0.48\textwidth]{pack-assessor-05.png}\hfill
\safeincludegraphics[width=0.48\textwidth]{pack-assessor-06.png}
\caption{Assessor pack, cards 1 to 6: preparation, session shape, the four moderation rules,
task scoring, the withheld success criteria, and the debrief questions.}
\end{figure}

\begin{figure}[H]
\centering
\safeincludegraphics[width=0.48\textwidth]{pack-assessor-07.png}\hfill
\safeincludegraphics[width=0.48\textwidth]{pack-assessor-08.png}
\[0.7em]
\safeincludegraphics[width=0.48\textwidth]{pack-assessor-09.png}\hfill
\hspace{0.48\textwidth}
\caption{Assessor pack, cards 7 to 9: the questionnaire structure, how the usability scale was
scored, and the troubleshooting sheet.}
\end{figure}
